% Abstract: most schools cap it at 350 words.
\chapter*{Abstract}
\addcontentsline{toc}{chapter}{Abstract}

This dissertation reconstructs the overland trade networks that linked the salt works of Lüneburg, the silver markets of the Saxon uplands, and the Hanseatic ports of the southern Baltic between 1350 and 1450. Historians have long treated the Hanse as a maritime league whose inland connections were secondary. Drawing on 4,118 entries from eleven toll registers, together with merchant correspondence and municipal account books, I argue instead that the land routes formed a dense, resilient network whose structure shaped the fortunes of the coastal cities.

The study proceeds in three steps. Chapter One sets out the historiographical problem and the questions the dissertation answers. Chapter Two describes the sources, the procedure used to reconcile divergent units of account, and a gravity model of trade that estimates the friction imposed by distance and tolls. Chapter Three applies these methods to the salt roads, showing that a handful of inland toll towns acted as brokers on the busiest paths between regional systems, and that the plague years of 1349--1351 and 1374 reorganised rather than destroyed this network.

The findings suggest that the prosperity of Lübeck rested as much on its control of overland brokerage as on its fleet, and that the much-discussed fifteenth-century decline of the Hanse began with a slow erosion of these inland positions.
\clearpage
